\dykband{EXERCISES IN SWAFL}
{\small Open SWAFL (desktop or online) with the benchmark project. The chapter of the \href{https://github.com/halilunaziru73-creator/Coupled-Water-Fluxes-Pore-Scale-Dynamics-Precision-Irrigation-Mulching-Feedback-Loops/raw/main/Software/v1.1.0/SWAFL_Desktop/docs/manual_v1.1.0/SWAFL_Manual_v1.1.0.pdf}{SWAFL manual} that explains each task is given in brackets; record your values in a table and compare them with what is expected (\emph{Check}).}\par\vspace{3pt}
\begin{exlist}
\item {\color{khblue}[Manual, Chapter 7]} Open \emph{Field balance} (the home layer) for the loam. Set \emph{Rainfall scale} to 0.6, 1.0 and 1.4 and run. Tabulate rain, irrigation, transpiration, evaporation and drainage, and check that each year closes.\par\noindent{\small\itshape Check: drainage rises much faster than transpiration as the rain rises.}
\item {\color{khblue}[Manual, Chapter 7]} In \emph{Profile evolution} press \emph{Play} and follow the wetting front after an irrigation. Estimate its depth after one and after three days.\par\noindent{\small\itshape Check: the front slows as it deepens, as the conductivity of the drier soil below is low.}
\item {\color{khblue}[Manual, Chapter 3]} Change the soil in the Properties panel from the loam to the sand and back. Note how the readily available water reported on the Dashboard changes.\par\noindent{\small\itshape Check: the sand holds a small fraction of the available water of the loam.}
\item {\color{khblue}[Manual, Chapter 14]} On the Dashboard open the \emph{Charts} tab and download the yearly water balance as a CSV file. Compute $\Delta S$ for one year by hand and compare with the program.\par\noindent{\small\itshape Check: the difference is below 1 mm.}
\end{exlist}
